
Compatible divisors were first introduced in \cite[Definition 2.7]{AZ_Kstab_fano_flag}, based on which the Abban-Zhuang method is developed to estimate (local) stability thresholds. In this section, we will use this concept to review the K-stability of toric varieties.

\begin{definition}
	A basis $\{s_1, \dots, s_{N_m}\}$ of $H^0(X, m L)$ is called \emph{compatible} with a valuation $v \in \DVal_X$ if for every $s \in H^0(X, m L)$,
	\begin{equation*}
		v(s) = \min_{c_j \neq 0} v(s_j),
	\end{equation*}
	where $s = \sum c_j s_j$, and $c_j \in \mathbb{C}$.
	% 
	When the basis $\{s_1, \dots, s_{N_m}\}$ is compatible with $v$, the $\mathbb{Q}$-divisor $D := {1 \over m N_m} \sum_j \div(s_j)$ is called an $m$-basis type $\mathbb{Q}$-divisor \emph{compatible} with $v$.
\end{definition}

The following lemma appears in the proof of \cite[Proposition 3.1]{AZ_Kstab_fano_flag} and is hidden in \cite[\S2.6 and Lemma 3.5]{BJ_thre_val_kst}. We write a short proof here for completeness.
\begin{lemma}\label[lemma]{lem:Smdef}
	For every valuation $v \in \DVal_X$, we have $S_m(L; v) = v(D)$ if the $m$-basis type $\mathbb{Q}$-divisor $D$ is compatible with $v$.
\end{lemma}
\begin{remark*}
	There is always some $m$-basis type $\mathbb{Q}$-divisor compatible with $v$. One can easily construct a compatible basis with the help of the valuative filtration, see \cite[\S3.1]{BJ_thre_val_kst}. Moreover, there always exists some compatible $\mathbb{Q}$-divisor compatible with any two valuations simultaneously, see \cite[Lemma 3.1]{AZ_Kstab_fano_flag}.
\end{remark*}
\begin{proof}
	Assume the basis $\{s_j\}_{j = 1}^{N_m}$ is compatible with $v$, and reorder it so that
	\begin{equation*}
		v(s_1) \leq \dots \leq v(s_{N_m}).
	\end{equation*}
	Given another basis $\{s_j'\}_{j = 1}^{N_m}$, similarly we can order it as above. Assume there exists a section $s_{j_0}'$ such that $v(s_{j_0}') > v(s_{j_0})$, then $v(s_{j_0 + 1}'), \dots, v(s_{N_m}') > v(s_{j_0})$ and we have a linearly independent subset
	\begin{equation*}
		\{s_{j_0}', s_{j_0 + 1}', \dots, s_{N_m}'\} \subseteq \{s \in H^0(X, m L) \mid v(s) > v(s_{j_0})\} \subseteq \<s_{j_0 + 1}, \dots, s_{N_m}\>_\mathbb{C},
	\end{equation*}
	which is impossible. Therefore, for every $j = 1, \dots, N_m$, we have $v(s_j') \leq v(s_j)$, and thus the basis $\{s_j\}_{j = 1}^{N_m}$ achieves the supremum.
\end{proof}


Recall that, to compute $S_m(L; v)$ by definition (see \Cref{def:S}), we need to consider all (and infinitely many) basis of $H^0(X, m L)$. The compatibility provides another way to compute $S_m(L; v)$. As valuations are invariant under (non-zero) scalar multiplication, we can generalize the compatibility from the bases to decompositions of vector spaces.
\begin{definition}\label[definition]{def:comp:decomp}
	For every valuation $v \in \DVal_X$, a decomposition $H^0(X, m L) =  \bigoplus_{J = 1}^k V_J$ (as a vector space) is \emph{compatible} with $v$ if 
	\begin{equation*}
		v(f) = \min_{J = 1, ..., k} v(f_J),
	\end{equation*}
	for every $f = \sum f_J \in \bigoplus_{J = 1}^k V_J$.
\end{definition}

Similar to compatible bases, we can use compatible decompositions to simplify the computation of $S_m(L; v)$: we only need to consider the basis of $H^0(X, mL)$ which is ``compatible'' with the decomposition.

\begin{lemma}\label[lemma]{lem:SmdefDecom}
	Under the assumption of \Cref{def:comp:decomp}, we have $S_m(L; v) = \sup_D v(D)$, where the $\mathbb{Q}$-divisor $D$ runs through all the $m$-basis type $\mathbb{Q}$-divisors defined from bases of the form $\bigcup_{J = 1}^k \mathcal{B}_J \subseteq H^0(X, m L)$ ($\mathcal{B}_J \subseteq V_J$ is a basis); equivalently,
	\begin{equation*}
		S_m(L; v) = {1 \over N_m} \sum_{J = 1}^k N_{m, J} \cdot S_m(V_J; v),
	\end{equation*}
	the (weighted) average, where $N_{m, J} := \dim(V_J)$.
\end{lemma}
\begin{remark*}
	When the subspace $V_J = \mathbb{C} s$, we have $S_m(V_J; v) = {1 \over m} v(s)$, this lemma reduces to \Cref{lem:Smdef}.
\end{remark*}
\begin{proof}
	For every subspace $V_J$, we can find a basis $\mathcal{B}_J \subseteq V_J$ compatible with $v$, and we have
	\begin{equation*}
		S_m(V_J; v) = {1 \over m N_{m, J}} \sum_j v(s_{J, j}),
	\end{equation*}
	then $\bigcup \mathcal{B}_J$ is a basis of $H^0(X, m L)$ compatible with $v$, which means
	\begin{equation*}
		S_m(L; v) = {1 \over m N_m} \sum_{J, j} v(s_{J, j}).
	\end{equation*}
\end{proof}


\subsection{Warm-up Application: Toric Varieties}\label{sec:comp:tor}
\newcommand{\TDValo}{\DVal^{T, \circ}}

In this section, we will find an anticanonical $\mathbb{Q}$-divisor computing the stability threshold of a toric variety. The idea is to prove that this $\mathbb{Q}$-divisor is compatible with all toric valuations. Then, in \Cref{sec:KstSph}, similarly, we will find an anticanonical $\mathbb{Q}$-divisor for spherical varieties.

We will follow the notations and some simple results from \cite[\S7]{BJ_thre_val_kst}, which is based on \cite{Fulton_intro_tor_vari}.

For a toric $T$-variety $X$, let $M := \ch(T)$, the character lattice, and $M_\mathbb{Q}: = M \otimes_\mathbb{Z} \mathbb{Q}$; let $\chi^u \in \mathbb{C}(T) = \mathbb{C}(X)$ be the character of $T$ associated to $u \in M$, thought of which as a rational function on $X$; let $X_1, \dots, X_k$ be the $T$-invariant prime divisors of $X$.

Given a line bundle $L$, we can pick a $T$-invariant divisor $D_0$ linearly equivalent to $L$, and let $P \subseteq M_\mathbb{Q}$ be the convex ($\mathbb{Q}$-)polytope associated to $D_0$.
% 
Now the subspace $H^0(X, m D_0) \subseteq \mathbb{C}(X)$ is a finite-dimensional $T$-submodule, and it can be decomposed into $1$-dimensional $T$-submodules, as
\begin{equation}
	H^0(X, m D_0) = \bigoplus_{u \in m P \cap M} \mathbb{C} \cdot \chi^u;
\end{equation}
we can define the $m$-basis type $\mathbb{Q}$-divisor
\begin{equation*}
	D_m = {1 \over m N_m} \sum_{u \in m P \cap M}( m D_0 + \div(\chi^u) ) = D_0 + {1 \over N_m} \sum {1 \over m} \div(\chi^u).
\end{equation*}
Now, let $D^T_L := D_0 + {1 \over m'} \div(\chi^{m' \bar{u}})$, where $\bar{u} \in P$ is the barycenter and $m' \bar{u} \in M$ for some $m' \in \mathbb{N} \setminus \{0\}$. The $\mathbb{Q}$-divisor $D^T_L$ is the limit of the sequence $(D_m)_m$ in the sense of coefficients; equivalently,
\begin{equation}\label{eq:toric:vDTL}
	v(D^T_L) = \lim_{m \to \infty} v(D_m)
\end{equation}
for every $T$-invariant valuation $v \in \DVal^T_X$.

The $\mathbb{Q}$-divisor $D^T_L$ is independent of the choice of $D_0$ (and thus the $T$-linearization). It is uniquely determined by the ample ($\mathbb{Q}$-)line bundle $L$. When $L$ is the anticanonical $\mathbb{Q}$-line bundle, we write $D^T_X$ for $D^T_{-\cano_X}$.

\begin{theorem}[c.f., {\cite[Theorem F and Corollary 7.4]{BJ_thre_val_kst}}]\label[theorem]{thm:toric}
	\begin{equation*}
	\delta(X) = \inf_{v \in \TDValo_X} {A_X(v) \over v(D^T_L)} = \min_{i = 1, \dots, k} {1 \over a_i},
	\end{equation*}
	where $a_1 X_1 + \dots + a_k X_k = D^T_X$, and $X_1, \dots, X_k$ are the $T$-invariant prime divisors.
\end{theorem}
\begin{remark*}
	By \cite[Proposition 7.3]{BJ_thre_val_kst}, the equality $\lct(D) = \inf_{v \in \TDValo_X} {A_X(v) \over v(D)}$ holds for every $T$-invariant $\mathbb{Q}$-divisor $D$, so we have $\delta(X) = \lct(D^T_X)$. No analogous result exists for computing the log canonical thresholds in the setting of spherical varieties. 

\end{remark*}

The formula of Blum and Jonsson for toric varieties is miraculously simple, and it completely solves the K-(semi)stability problem of toric varieties. We are going to reprove this theorem in two steps. The first step is to use the compatible divisors to prove,
\begin{lemma}\label[lemma]{lem:toric:comp}
	For every $v \in \TDValo_X$ and all sufficiently divisible integer $m$, we have $S(v) = v(D^T_L)$, and thus $\delta_T(L) = \inf_{v \in \TDValo_X} {A_X(v) \over v(D^T_L)}$ (see \Cref{sec:prelim:kstab:GKst}).
\end{lemma}
\begin{proof}
	We only need to prove for every $v \in \DVal^T_X$, the equality $S_m(v) = v(D_m)$ holds, i.e., to prove $D_m$ is compatible with $v$. However, $T$-invariant valuations are exactly ``toric valuations'', which means,
	\begin{equation*}
		v(\sum_{u \in M} c_u \chi^u) = \min_{\substack{u \in M \\ c_u \neq 0}} v(\chi^u).
	\end{equation*}
	This is exactly the definition of compatibility.
\end{proof}

An intermediate result is
\begin{equation}\label{eq:deltaTm}
	\delta_{m, T}(L) = \inf_{v \in \TDValo_X} {A_X(v) \over v(D_m)},
\end{equation}
which we will use to prove the next lemma.

The second step is
\begin{lemma}[c.f., {\cite[Corollary 7.15]{BJ_thre_val_kst}}]\label[lemma]{lem:toric:delta}
	$\delta(X) = \delta_T(X)$.
\end{lemma}
\begin{proof}
	The original proof of Blum and Jonsson deforms the general valuative filtrations into their initial filtrations (see \cite[\S7.4]{BJ_thre_val_kst}). We are going to use Zhuang's theorem (see \Cref{thm:Zhuang}) instead.

	The standard anticanonical divisor is $X_1 + \dots + X_k$ for the toric variety $X$, and the effective $\mathbb{Q}$-divisor $D_m$ is of the form $D_m = a_1 X_1 + \dots + a_k X_k$ for sufficiently divisible integer $m$. There must be an element $a_i \geq 1$ (otherwise, we have $D_m < X_1 + \dots + X_k$), and by \Cref{eq:deltaTm},
	\begin{equation*}
		\delta_{m, T}(X) \leq {A_X(\ord_{X_i}) \over \ord_{X_i}(D_m)} = {1 \over a_i} \leq 1.
	\end{equation*}
	After taking the limit we have $\delta_T(X) \leq 1$. As $\delta(X) \leq \delta_T(X) \leq 1$, we can apply Zhuang's theorem to show $\delta(X) = \delta_T(X)$.
\end{proof}

\begin{proof}[Proof of \Cref{thm:toric}]
	We just combine \Cref{lem:toric:comp} and \Cref{lem:toric:delta}. The last equality is a well known property of toric varieties, see, for example, \cite[Proposition 7.2]{BJ_thre_val_kst}.
\end{proof}

Recall that, a $\mathbb{Q}$-Fano variety $X$ is K-semistable iff $\delta(X) \geq 1$, we have
\begin{corollary}\label[corollary]{cor:tor:Kss}
	A $\mathbb{Q}$-Fano toric variety $X$ is K-semistable iff the $\mathbb{Q}$-divisor $D^T_X = X_1 + \dots + X_k$.
\end{corollary}
\begin{remark*}
	This is just another way of expressing \cite[Corollary 7.17]{BJ_thre_val_kst}.
\end{remark*}

Spherical varieties are a generalization of toric varieties. We are going to construct similar $\mathbb{Q}$-divisors for $\mathbb{Q}$-Fano spherical varieties.
